\documentclass[11pt,a4paper]{article}
\usepackage[margin=2cm]{geometry}
\usepackage[T1]{fontenc}
\usepackage{lmodern}
\usepackage{amsmath}
\usepackage{booktabs}
\usepackage{tabularx}
\usepackage{enumitem}
\usepackage{xcolor}
\usepackage[colorlinks=true,linkcolor=blue!50!black,urlcolor=blue!50!black]{hyperref}
\setlist{nosep,leftmargin=*}
\setlength{\parskip}{4pt}
\setlength{\parindent}{0pt}
\newcommand{\tbd}[1]{\textcolor{red}{[TBD: #1]}}
\newcommand{\teamcontacts}{}
\IfFileExists{team_private.tex}{\input{team_private.tex}}{\IfFileExists{docs/team_private.tex}{\input{docs/team_private.tex}}{}}

\title{\textbf{Industrial Defect Segmentation on MVTec AD}\\[2pt]\large Project Plan --- Computer Vision 2026}
\author{jestersw (team lead) \and moonshine1l \and psandersi \and AigulFar}
\date{6 October 2026\\[2pt]\small Repository: \url{https://github.com/jestersw/Defect-Segmentation}}

\begin{document}
\maketitle

\section{Project Scope}

\textbf{Objective.} We segment surface defects such as cracks, scratches, holes and other damaged regions on industrial objects from the MVTec AD dataset, producing a pixel-level defect mask for every image.

\textbf{Research question.} \emph{Does replacing the binary cross-entropy (BCE) loss with a combined Dice + Focal loss improve the recall of small defects in a U-Net segmentation model, without a large drop in precision?}

\textbf{Motivation and hypothesis.} Defects usually occupy a small fraction of the image, so the BCE loss is dominated by the many easy background pixels and tiny defects contribute little to the gradient. The Dice loss measures overlap and is insensitive to the foreground/background imbalance, while the Focal loss down-weights easy pixels. We therefore expect the combined loss to recover more small defects, at the risk of more false-positive pixels.

\textbf{Exact setting.}
\begin{itemize}
  \item Task: binary semantic segmentation, defect (1) vs.\ background (0); all defect types of a category are merged into one class.
  \item Categories: \texttt{tile}, \texttt{hazelnut}, \texttt{wood}, \texttt{metal\_nut}, \texttt{capsule}.
  \item Model: U-Net with a ResNet34 encoder pretrained on ImageNet.
  \item Compared conditions: baseline (BCE) vs.\ condition B (Dice + Focal); the loss is the only changed factor.
  \item Ablation: input resolution $512$ vs.\ the default $256$ for the better condition, plus a threshold sensitivity curve.
\end{itemize}

\textbf{Out of scope.} Unsupervised anomaly detection without masks, multi-class defect-type segmentation, and the design of new architectures.

\section{Dataset Preparation}

\textbf{Dataset.} MVTec AD~\cite{mvtec} contains high-resolution images of 15 object and texture categories with pixel-precise ground-truth masks for defective images. It is distributed under the CC BY-NC-SA 4.0 license and is used here for non-commercial educational purposes only. We use five categories that contain cracks, scratches or other surface damage (Table~\ref{tab:data}). The data is obtained from the Kaggle mirror \url{https://www.kaggle.com/datasets/ipythonx/mvtec-ad}, preserving the original folder structure. The downloaded categories contain 1,860 images (1,307 official training good, 137 test good, and 416 test defective) with 416 masks. The script \texttt{scripts/dataset\_stats.py} verifies the expected inventory, readability, matching image/mask dimensions, and absence of missing, orphaned or empty defect masks; all checks pass. Original counts are recorded in \texttt{results/dataset\_counts.csv}, before the exclusion below.

\textbf{Why we re-split.} The official MVTec AD training set contains only defect-free images, and masks exist only in the official test set. Since we train a supervised segmentation model, we build our own split from the official data:
\begin{itemize}
  \item defective images of the official test set of the five categories, except \texttt{metal\_nut/flip};
  \item good images: the official \texttt{test/good} images plus a random sample of \texttt{train/good}, so that the number of good images equals the number of defective images per category. Good images teach the model not to predict defects on clean surfaces.
\end{itemize}
The defect type \texttt{metal\_nut/flip} (23 images) is excluded: its masks cover about 48\% of the image, i.e.\ the whole object, so it represents a global orientation anomaly rather than a localized surface defect and could dominate overlap metrics. The exclusion is configured in \texttt{data.exclude\_defect\_types} and is applied before balancing good images and generating splits. Inventory checks and descriptive statistics still cover all original masks; the training threshold uses the filtered split. The resulting pool contains 393 defective images and 393 good images (all 137 official test good images plus 256 sampled official training good images).

\begin{table}[h]
\centering
\small
\setlength{\tabcolsep}{4pt}
\begin{tabular}{lrrrrrrr}
\toprule
 & \multicolumn{3}{c}{Official MVTec AD} & & \multicolumn{3}{c}{Our split (70/15/15)} \\
\cmidrule(lr){2-4}\cmidrule(lr){6-8}
Category & Train good & Test good & Test defect & Pool & Train & Val & Test \\
\midrule
capsule    & 219 & 23 & 109 & 218 & 152 & 34 & 32 \\
hazelnut   & 391 & 40 & 70  & 140 & 99  & 20 & 21 \\
metal\_nut & 220 & 22 & 93 (70$^{*}$) & 140 & 97 & 21 & 22 \\
tile       & 230 & 33 & 84  & 168 & 117 & 25 & 26 \\
wood       & 247 & 19 & 60  & 120 & 85  & 18 & 17 \\
\midrule
Total      & 1307 & 137 & 416 (393) & 786 & 550 & 118 & 118 \\
\bottomrule
\end{tabular}
\caption{Official MVTec AD counts~\cite{mvtec} (verified on the downloaded data) and our split. Pool = used defective images + an equal number of good images. $^{*}$23 \texttt{flip} images excluded. Split counts are produced by \texttt{scripts/make\_split.py} and stored in \texttt{results/split\_summary.csv}.}
\label{tab:data}
\end{table}

\textbf{Split procedure.} The pool is split $70/15/15$ into training, validation and test sets using two stratified \texttt{train\_test\_split} calls with seed $42$ (first $70/30$, then $50/50$ of the held-out part). Strata are (category, defect type), with one good-image stratum per category. The five categories contain 23 defect types (22 after the exclusion); every retained type appears in all three generated splits. Every stratum has at least eight images. The fixed manifests \path{splits/train.txt}, \path{splits/val.txt}, and \path{splits/test.txt} are reused in all experiments.

\textbf{Leakage prevention.}
\begin{itemize}
  \item Each pooled image path appears in exactly one split; byte-identical files across splits are rejected by an MD5 check. These checks pass; near-duplicate images and acquisition-group overlap have not been assessed.
  \item Under the planned evaluation protocol, the test split will be used only for final evaluation; checkpoint selection, early stopping and probability-threshold tuning will use validation data.
  \item Planned image normalisation uses ImageNet statistics. The small-defect area threshold is estimated on training masks only.
  \item The split script refuses to overwrite committed split files unless explicitly forced, and SHA-256 hashes of the split files are stored in \texttt{results/split\_metadata.json} and checked in CI.
\end{itemize}

\textbf{Preprocessing.} Stored images retain their original resolution. The implemented statistics pipeline resizes masks to $512\times512$ using OpenCV nearest-neighbour interpolation and binarises them at $>0$. The planned training loader will resize images with bilinear interpolation and masks with nearest-neighbour interpolation to $256\times256$ ($512\times512$ in the ablation), and normalise images with ImageNet mean and standard deviation. These training transforms are planned; \texttt{src/data.py} is not yet implemented.

\textbf{Augmentation (train only, planned).} Horizontal and vertical flips and $90^\circ$ rotations will be applied jointly to images and masks. Mild brightness and contrast changes will affect images only. Validation and test will use deterministic preprocessing without augmentation.

\textbf{Defect-size statistics.} Across all 416 original masks, the descriptive statistics contain 638 connected components at $512\times512$, extracted with 8-connectivity. Areas range from 1 to 127,709 pixels; no defect mask becomes empty after resizing. A component is called \emph{small} if its area is below the 33rd percentile of training component areas: $859.28$ pixels, approximately $0.3278\%$ of the image. This threshold is computed with NumPy linear interpolation from the 449 components in the 550 training images only, with no contribution from validation or test masks. The full-precision value and split hashes are saved in \texttt{results/small\_defect\_threshold.json}; the same value is used in \texttt{configs/base.yaml}. The exclusion rule is recorded in \texttt{results/split\_metadata.json}. The CSV and histogram describe all original defect types, including \texttt{flip}, whereas the threshold uses training components exclusively and excludes \texttt{flip}.

\section{Experimental Plan}

\subsection{Compared conditions}

\begin{table}[h]
\centering
\small
\begin{tabularx}{\linewidth}{lXX}
\toprule
 & Baseline & Condition B \\
\midrule
Model & U-Net, ResNet34 encoder, ImageNet weights & same \\
Loss & $\mathcal{L}_{\mathrm{BCE}}$ & $0.5\,\mathcal{L}_{\mathrm{Dice}} + 0.5\,\mathcal{L}_{\mathrm{Focal}}$ ($\gamma=2$, $\alpha=0.25$) \\
Data, resolution, augmentation & split of Section 2, $256\times256$ & same \\
Optimisation & AdamW, lr $10^{-4}$, wd $10^{-4}$, cosine schedule, batch 8, $\le 60$ epochs & same \\
Model selection & early stopping on validation Dice, patience 10 & same \\
Threshold & $0.5$ & same \\
\bottomrule
\end{tabularx}
\caption{Experimental conditions. Only the loss differs.}
\end{table}

With predicted probabilities $p_i$ and ground truth $g_i\in\{0,1\}$ over pixels $i$:
\begin{align*}
\mathcal{L}_{\mathrm{BCE}} &= -\frac{1}{N}\sum_i \big[g_i\log p_i + (1-g_i)\log(1-p_i)\big],\\
\mathcal{L}_{\mathrm{Dice}} &= 1 - \frac{2\sum_i p_i g_i + \varepsilon}{\sum_i p_i + \sum_i g_i + \varepsilon},\qquad
\mathcal{L}_{\mathrm{Focal}} = -\frac{1}{N}\sum_i \alpha_t\,(1-p_{t,i})^{\gamma}\log p_{t,i},
\end{align*}
where $p_{t,i}=p_i$ if $g_i=1$ and $1-p_i$ otherwise. BCE weights every pixel equally, Dice depends only on the overlap with the defect, and Focal reduces the weight of pixels that are already classified confidently.

\textbf{Training strategy.} Both conditions fine-tune the full network (encoder and decoder) starting from ImageNet weights, with identical schedule and hyperparameters. Because only CPU compute is available (Section 4), each condition is trained with a single seed (42); predictions are upsampled to $512\times512$ before evaluation so that the small-defect threshold from Section 2 applies unchanged.

\textbf{Ablation and sensitivity.}
\begin{itemize}
  \item Resolution: the better condition is retrained at $512\times512$. Hypothesis: the higher input resolution preserves small defects and improves their recall.
  \item Threshold: masks are binarised at $0.3$, $0.5$ and $0.7$ to show the precision--recall trade-off.
\end{itemize}

\textbf{Decision rule.} Condition B answers the research question positively if small-defect recall improves over the baseline while pixel precision and the good-image false-positive rate do not degrade substantially.

\subsection{Evaluation metrics}

Let $P$ and $G$ be the predicted and ground-truth defect pixel sets of an image.
\begin{itemize}
  \item \textbf{Dice} $=\dfrac{2|P\cap G|}{|P|+|G|}$ and \textbf{IoU} $=\dfrac{|P\cap G|}{|P\cup G|}$, averaged over defective test images. They measure mask quality; on good images $G=\emptyset$, so these metrics are undefined there and good images are evaluated by the false-positive rate instead.
  \item \textbf{Pixel precision} $=\dfrac{TP}{TP+FP}$ and \textbf{pixel recall} $=\dfrac{TP}{TP+FN}$, aggregated over all test pixels. They show which kind of error dominates.
  \item \textbf{Small-defect recall}: share of small ground-truth components (Section 2) of which at least $25\%$ of the pixels are predicted as defect. This metric answers the research question directly, because large defects dominate Dice and pixel recall.
  \item \textbf{Good-image false-positive rate (FPR)}: share of good test images where the prediction contains a connected component of at least 20 pixels. It guards against a model that wins recall by predicting defects everywhere.
  \item \textbf{Runtime}: CPU training time per epoch, inference time per image (after warm-up), and peak process RAM (the selected compute is described in Section 4).
\end{itemize}
Results are reported overall and per category in one main comparison table, together with the ablation table, the threshold curve and the runtime comparison.

\section{Implementation Plan}

\textbf{Pipeline.}
Data preparation, split checks, metrics and the standalone smoke test are implemented. The dataset loader, model wrapper, losses, training, evaluation and profiling modules in \texttt{src/} are currently placeholders; the training and evaluation steps below specify the implementation contract.
\begin{enumerate}
  \item Download the five categories and verify the image/mask inventory with \path{scripts/download_data.py} and \path{scripts/dataset_stats.py}.
  \item Build the stratified split with \path{scripts/make_split.py}, run leakage checks, then compute the train-only small-defect threshold using \path{scripts/dataset_stats.py} with \texttt{--splits-dir splits}.
  \item Verify U-Net forward/backward and measure synthetic training steps with \path{scripts/smoke_test.py} at both resolutions.
  \item PyTorch \texttt{Dataset}/\allowbreak\texttt{DataLoader} with preprocessing and train-only augmentation.
  \item Build the U-Net model and train it with the configured loss; write per-epoch CSV logs and learning curves; keep the checkpoint with the best validation Dice.
  \item Evaluate the best checkpoint on the test split and save metrics as JSON.
  \item Export predictions for failure analysis and run the ablation and runtime measurements.
\end{enumerate}

\textbf{Repository structure.}
\begin{center}
\small
\setlength{\tabcolsep}{4pt}
\begin{tabularx}{\linewidth}{lX}
\toprule
File & Responsibility \\
\midrule
\texttt{scripts/download\_data.py}, \texttt{scripts/dataset\_stats.py} & data acquisition and statistics \\
\texttt{scripts/make\_split.py}, \texttt{splits/*.txt} & fixed split and leakage checks \\
\texttt{scripts/smoke\_test.py} & CPU step timing and peak process RAM measurement \\
\texttt{src/data.py} & planned dataset, preprocessing, augmentation \\
\texttt{src/model.py}, \texttt{src/train.py} & planned model wrapper and training loop \\
\texttt{src/losses.py} & planned BCE and Dice + Focal losses \\
\texttt{src/metrics.py}, \texttt{src/evaluate.py} & implemented metrics; planned evaluation loop \\
\texttt{src/profiling.py} & planned end-to-end training/inference profiling \\
\texttt{configs/base.yaml} & all hyperparameters \\
\texttt{requirements.txt}, \texttt{.python-version} & pinned libraries and Python 3.11 \\
\texttt{tests/} & unit tests (metrics, config, statistics, split builder, leakage) \\
\texttt{.github/workflows/ci.yml} & CI: lint, tests and CPU check on main/s1 pushes and PRs \\
\bottomrule
\end{tabularx}
\end{center}

\textbf{Planned interfaces.} A dataset item is \texttt{(image, mask)} with shapes $[3,H,W]$ and $[1,H,W]$, mask values in $\{0,1\}$; the model returns logits of shape $[B,1,H,W]$; \texttt{evaluate(model, loader, threshold)} returns a dictionary with \texttt{dice}, \texttt{iou}, \texttt{precision}, \texttt{recall}, \texttt{small\_recall} and \texttt{good\_fpr}. Training hyperparameters will be read from \texttt{configs/base.yaml}; condition B differs from the baseline config only in the \texttt{loss} field.

\textbf{Training loop (planned).} Fully fine-tune encoder and decoder using AdamW, learning rate $10^{-4}$, weight decay $10^{-4}$, cosine annealing over at most 60 epochs and batch size 8. Select the best checkpoint by maximum validation Dice; stop after 10 consecutive epochs without improvement. Save per-epoch training/validation loss, validation Dice and learning rate to CSV, and plot loss and Dice curves. Baseline uses \texttt{BCEWithLogitsLoss}; condition B uses the configured Dice + Focal loss.

\textbf{Libraries and resources.} Python 3.11; PyTorch 2.6.0 and torchvision 0.21.0; segmentation-models-pytorch 0.4.0 (U-Net, ResNet34 ImageNet weights, Dice and Focal losses); Albumentations 2.0.8; opencv-python and its Albumentations dependency opencv-python-headless 4.11.0.86; scikit-image 0.25.2 (connected components); NumPy 2.1.3; scikit-learn 1.9.0; PyYAML 6.0.2; Matplotlib 3.10.3; tqdm 4.67.1; psutil 6.1.1 (process RAM). These direct dependencies are pinned in \path{requirements.txt}; CI tools are pinned separately in \path{requirements-dev.txt}.

\textbf{Reproducibility.} The smoke test seeds Python, NumPy and PyTorch with 42, enables deterministic algorithms, fixes CPU thread counts, and saves its configuration, package versions, device information and train-manifest hash with each JSON result. Deterministic cuDNN is also set with benchmarking and TF32 disabled, but cuDNN is inactive for CPU execution. Training will use the same settings, seeded DataLoader generators/workers, and a resolved configuration saved per run next to the checkpoint, CSV and curves. One seed per condition is used because of the CPU-only budget. Exact reproduction is scoped to the same software and hardware environment. Datasets and checkpoints are not stored in git.

\textbf{Continuous integration.} GitHub Actions runs on pushes to \texttt{main} and \texttt{s1/**}, and on pull requests. It installs pinned CPU PyTorch and project requirements on Python 3.11, checks dependency consistency, lints with Ruff, runs pytest, and checks a small U-Net forward/backward pass on CPU without downloading pretrained weights. Unit tests use synthetic data and cover metrics, configuration, statistics, the split builder, leakage, committed manifest hashes and runtime projection arithmetic. The small CI check does not replace the full local benchmark. Project policy requires passing CI and one approval before merging into \texttt{main}.

\textbf{Compute.} We use local CPU compute: a laptop with an AMD Ryzen AI 9 HX 370 (12 cores) and 31.12\,GiB RAM, running Windows and Python 3.11.12; PyTorch uses 8 computation threads. \path{scripts/smoke_test.py} ran 20 FP32 forward/backward/AdamW steps after 3 warm-up steps with batch size 8 and ImageNet weights; the output shapes were $[8,1,512,512]$ and $[8,1,256,256]$. Raw results are in \path{results/smoke_test_cpu.json}; the exact command is in \path{docs/compute.md}.

\begin{center}
\small
\begin{tabular}{rrrrr}
\toprule
Resolution & Mean s/step & Peak RAM (GiB) & Min/epoch (69 steps) & Hours/60 epochs \\
\midrule
$512^2$ & 7.60 & 4.81 & 8.74 & 8.74 \\
$256^2$ & 1.91 & 2.25 & 2.19 & 2.19 \\
\bottomrule
\end{tabular}
\end{center}

These projections cover training steps only and exclude data loading, augmentation, validation and checkpoint I/O. One 60-epoch run at $512^2$ therefore takes about 9\,h. We train both conditions at $256^2$ (at most about 2.2\,h per run) and use $512^2$ only for the single ablation run; early stopping is expected to shorten runs.

\section{Team Responsibilities}

\begin{table}[h]
\centering
\small
\begin{tabularx}{\linewidth}{lX}
\toprule
Member & Responsibility \\
\midrule
jestersw (lead) & Research question and experimental design; metric definitions and evaluation module; comparison tables; MLOps: CI pipeline with automated tests and branch protection; repository integration and reproducibility; reports and demo. \\
moonshine1l & Dataset acquisition and category selection; stratified split with leakage checks; preprocessing and augmentation; defect-size statistics. \\
psandersi & Baseline training pipeline (U-Net + BCE): training loop, checkpoints, curve logging; CPU compute measurement. \\
AigulFar & Condition B (Dice + Focal) under identical settings; resolution ablation and threshold sensitivity; runtime and memory comparison. \\
All members & At least one analysed failure case each; peer review of pull requests. \\
\bottomrule
\end{tabularx}
\end{table}

\teamcontacts

\textbf{Workflow.} All changes go through pull requests that require a passing CI run and one approval. Review pairs are moonshine1l $\leftrightarrow$ jestersw and psandersi $\leftrightarrow$ AigulFar, so every member can explain a neighbouring part of the project.

\begin{table}[h]
\centering
\small
\begin{tabularx}{\linewidth}{lXl}
\toprule
Date & Stage & Main owners \\
\midrule
13 Oct & First working model: dataset code, baseline results & moonshine1l, psandersi, jestersw \\
20 Oct & Second approach: Dice + Focal working & AigulFar \\
27 Oct & Progress check: preliminary comparison, remaining experiments & jestersw, AigulFar \\
3 Nov & Results: final metrics, runtime, $\ge 3$ failure cases & all \\
10 Nov & Final submission and demo & jestersw, all \\
\bottomrule
\end{tabularx}
\end{table}

\begin{thebibliography}{9}
\bibitem{mvtec} P.~Bergmann, M.~Fauser, D.~Sattlegger, C.~Steger. MVTec AD --- A Comprehensive Real-World Dataset for Unsupervised Anomaly Detection. \emph{CVPR}, 2019.
\end{thebibliography}

\end{document}
